\section*{अध्याय \ref{ch:b1:quantum-numbers} --- क्वांटम संख्याएँ और इलेक्ट्रॉनिक विन्यास}

\begin{solution}{exo:b1:quantum-numbers:1}
(a) असंभव: $l \le n - 1 = 1$। (b) संभव (एक 3p इलेक्ट्रॉन)।
(c) असंभव: $m_s = \pm\tfrac12$। (d) संभव (एक 4f इलेक्ट्रॉन)।
(e) असंभव: $l = 0$ के लिए $m_l$ केवल 0 हो सकता है।
\end{solution}

\begin{solution}{exo:b1:quantum-numbers:2}
3d: 5 कक्षक, 10 इलेक्ट्रॉन; 4f: 7 कक्षक, 14 इलेक्ट्रॉन; 5p: 3
कक्षक, 6 इलेक्ट्रॉन; कोश $n = 2$: $n^2 = 4$ कक्षक, $2n^2 = 8$
इलेक्ट्रॉन।
\end{solution}

\begin{solution}{exo:b1:quantum-numbers:3}
\ce{O}: $1\text{s}^2 2\text{s}^2 2\text{p}^4 = [\ce{He}]\,2\text{s}^2 2\text{p}^4$,
6 \omterm{def:b1:quantum-numbers:configuration}{संयोजकता इलेक्ट्रॉन}। \ce{P}: $1\text{s}^2 2\text{s}^2 2\text{p}^6 3\text{s}^2
3\text{p}^3 = [\ce{Ne}]\,3\text{s}^2 3\text{p}^3$, 5। \ce{K}: $[\ce{Ar}]\,4\text{s}^1$,
1। \ce{Ca}: $[\ce{Ar}]\,4\text{s}^2$, 2। \ce{Br}: $[\ce{Ar}]\,3\text{d}^{10}
4\text{s}^2 4\text{p}^5$, 7 (भरा हुआ 3d क्रोड का भाग है)।
% ledger: gs:K, gs:Ca
\end{solution}

\begin{solution}{exo:b1:quantum-numbers:4}
\ce{N}, 2p: \omorbs{u,u,u}, 3 अयुग्मित। \ce{S}, 3p: \omorbs{ud,u,u}, 2
अयुग्मित। \ce{Cl}, 3p: \omorbs{ud,ud,u}, 1 अयुग्मित। हर स्थिति में 2s या
3s पूरा भरा है, \omorbs{ud}।
\end{solution}

\begin{solution}{exo:b1:quantum-numbers:5}
\ce{Na+}, \ce{Mg^{2+}}, \ce{Al^{3+}}, \ce{F-}, \ce{O^{2-}}: $1\text{s}^2
2\text{s}^2 2\text{p}^6$, नियॉन के समइलेक्ट्रॉनी। \ce{Cl-}, \ce{S^{2-}},
\ce{Ca^{2+}}: $[\ce{Ne}]\,3\text{s}^2 3\text{p}^6$, आर्गन के
समइलेक्ट्रॉनी।
\end{solution}

\begin{solution}{exo:b1:quantum-numbers:6}
4s इलेक्ट्रॉन पहले निकलते हैं (\cref{prop:b1:quantum-numbers:ions})।
\ce{Ti^{2+}} के लिए $[\ce{Ar}]\,3\text{d}^2$, 2 अयुग्मित; \ce{Mn^{2+}} के लिए
$3\text{d}^5$, 5; \ce{Co^{2+}} के लिए $3\text{d}^7$, 3 (\omorbs{ud,ud,u,u,u});
\ce{Ni^{2+}} के लिए $3\text{d}^8$, 2; \ce{Zn^{2+}} के लिए $3\text{d}^{10}$, 0।
% ledger: gs:Ti2+, gs:Mn2+, gs:Co2+, gs:Zn2+
\end{solution}

\begin{solution}{exo:b1:quantum-numbers:7}
$\Delta E = 13.6\,(1/9 - 1/16) = 13.6 \times 7/144 = \qty{0.661}{eV}$;
$\lambda = 1240/0.661 = \qty{1876}{nm} \approx \qty{1.88}{\micro m}$:
अवरक्त (पाशन श्रेणी की पहली रेखा)।
\end{solution}

\begin{solution}{exo:b1:quantum-numbers:8}
(a) $[\ce{Ne}]\,3\text{s}^2 3\text{p}^4$: सल्फ़र, $Z = 16$।
(b) $[\ce{Ar}]\,3\text{d}^{10} 4\text{s}^2 4\text{p}^5$: ब्रोमीन, $Z = 35$।
(c) $[\ce{Kr}]\,5\text{s}^1$: रुबिडियम, $Z = 37$।
(d) $[\ce{Ar}]\,3\text{d}^3 4\text{s}^2$: वैनेडियम, $Z = 23$।
सल्फ़र का एक 3p इलेक्ट्रॉन: उदाहरण के लिए $(3, 1, -1, +\tfrac12)$।
% ledger: gs:V
\end{solution}

\begin{solution}{exo:b1:quantum-numbers:9}
4s के लिए $n + l = 4 + 0 = 4$; 3d के लिए $n + l = 3 + 2 = 5$: 4s पहले
आता है। 3p ($n + l = 4$) के बाद अगला \omterm{def:b1:quantum-numbers:orbital}{उपकोश} 4s है ($n + l = 4$, बड़ा
$n$), जो आवर्त 4 को खोलता है; 3d ($n + l = 5$) 4s के बाद ही आता है।
इसलिए आवर्त 3 में केवल 3s और 3p होते हैं: $2 + 6 = 8$ तत्व।
\end{solution}

\begin{solution}{exo:b1:quantum-numbers:10}
आयनन ऊर्जा $Z^2E_R = 4 \times 13.6 = \qty{54.4}{eV}$, हाइड्रोजन की
चार गुनी। $2 \to 1$: $\Delta E = 54.4\,(1 - 1/4) = \qty{40.8}{eV}$,
$\lambda = 1240/40.8 = \qty{30.4}{nm}$, \qty{121.6}{nm} वाली हाइड्रोजन
रेखा से चार गुनी छोटी: हर ऊर्जा $Z^2 = 4$ से गुणा हो जाती है।
\end{solution}

\begin{solution}{exo:b1:quantum-numbers:11}
(a) बेरिलियम की \omterm{def:b1:quantum-numbers:energy-level}{उत्तेजित अवस्था} (\omterm{def:b1:quantum-numbers:energy-level}{मूल अवस्था} $1\text{s}^2 2\text{s}^2$)।
(b) असंभव: 2p में अधिक-से-अधिक 6 इलेक्ट्रॉन आते हैं। (c) बेरिलियम की
\omterm{def:b1:quantum-numbers:energy-level}{उत्तेजित अवस्था}। (d) फ़ॉस्फ़ोरस की \omterm{def:b1:quantum-numbers:energy-level}{मूल अवस्था}। (e) असंभव: s \omterm{def:b1:quantum-numbers:orbital}{उपकोश} में
दो ही इलेक्ट्रॉन आते हैं (पाउली)।
\end{solution}

\begin{solution}{exo:b1:quantum-numbers:12}
\ce{Fe^{3+}} $3\text{d}^5$: 5; \ce{Mn^{2+}} $3\text{d}^5$: 5; \ce{Fe^{2+}}
$3\text{d}^6$: 4; \ce{Cu^{2+}} $3\text{d}^9$: 1; \ce{Zn^{2+}}
$3\text{d}^{10}$: 0। क्रम: $\ce{Fe^{3+}} = \ce{Mn^{2+}} > \ce{Fe^{2+}}
> \ce{Cu^{2+}} > \ce{Zn^{2+}}$। आयरन(III) में आयरन(II) से एक \omterm{def:b1:quantum-numbers:unpaired}{अयुग्मित
इलेक्ट्रॉन} अधिक है, इसलिए वह अधिक आकर्षित होता है; ज़िंक(II) में कोई नहीं,
इसलिए वह आकर्षित नहीं होता।
% ledger: gs:Fe2+, gs:Fe3+, gs:Mn2+, gs:Cu2+, gs:Zn2+
\end{solution}

\begin{solution}{pb:b1:quantum-numbers:1}
\textbf{1.} $E_1 = \qty{-13.6}{eV}$, $E_2 = \qty{-3.40}{eV}$,
$E_3 = \qty{-1.51}{eV}$, $E_4 = \qty{-0.85}{eV}$।
\textbf{2.} $\Delta E = 3.40 - 1.51 = \qty{1.89}{eV}$, $\lambda = 1240/1.89
= \qty{656}{nm}$: लाल।
\textbf{3.} $4 \to 2$: \qty{2.55}{eV}, \qty{486}{nm} (नीला-हरा);
$5 \to 2$: \qty{2.86}{eV}, \qty{434}{nm} (बैंगनी); $6 \to 2$: \qty{3.022}{eV},
\qty{410}{nm} (बैंगनी)।
\textbf{4.} $7 \to 2$: \qty{3.12}{eV}, \qty{397}{nm}, पहले ही पराबैंगनी;
ऊँचे संक्रमण और भी छोटे हैं, और $3 \to 2$ बामर की सबसे लंबी
तरंगदैर्घ्य है। लाइमन रेखाएँ सब पराबैंगनी में हैं और पाशन
रेखाएँ सब अवरक्त में, इसलिए केवल चार रेखाएँ दिखती हैं।
\textbf{5.} $\infty \to 2$: \qty{3.40}{eV}, $\lambda = \qty{365}{nm}$।
\textbf{6.} $\Delta E = 13.6 \times 3/4 = \qty{10.2}{eV}$, $\lambda =
\qty{122}{nm}$: सुदूर पराबैंगनी, आँख को अदृश्य और काँच तथा वायु द्वारा
अवशोषित।
\textbf{7.} फ़ोटॉन को कम-से-कम \qty{13.6}{eV} लाना चाहिए: $\lambda \le
1240/13.6 = \qty{91.2}{nm}$।
\textbf{8.} $l = 0$ ($m_l = 0$), $l = 1$ ($m_l = -1, 0, 1$), $l = 2$
($m_l = -2, \dots, 2$): $1 + 3 + 5 = 9$ कक्षक।
\textbf{9.} $\sum_{l=0}^{n-1}(2l+1) = n^2$ कक्षक, हर एक में दो इलेक्ट्रॉन:
$2n^2$; $n = 4$ के लिए 32।
\textbf{10.} 1s, 2s, 2p, 3s, 3p, 4s, 3d, 4p, 5s, 4d, 5p, 6s, 4f, 5d, 6p,
7s, 5f, 6d, 7p।
\textbf{11.} हर आवर्त $n$s से $n$p तक चलता है: 2, 8, 8, 18, 18, 32, 32।
\textbf{12.} 3d ($n + l = 5$) 4s ($n + l = 4$) के बाद आता है, और 4s
आवर्त 4 को खोलता है।
\textbf{13.} $Z = 2, 10, 18, 36, 54, 86, 118$।
\textbf{14.} $[\ce{Ar}]\,3\text{d}^6 4\text{s}^2$; 3d: \omorbs{ud,u,u,u,u},
4 \omterm{def:b1:quantum-numbers:unpaired}{अयुग्मित इलेक्ट्रॉन}।
\textbf{15.} \ce{Fe^{2+}} $[\ce{Ar}]\,3\text{d}^6$: 4 अयुग्मित;
\ce{Fe^{3+}} $[\ce{Ar}]\,3\text{d}^5$: 5 अयुग्मित, आधा भरा \omterm{def:b1:quantum-numbers:orbital}{उपकोश}।
\textbf{16.} अनुमानित $3\text{d}^4 4\text{s}^2$: 4 अयुग्मित; प्रेक्षित
$3\text{d}^5 4\text{s}^1$: 6 अयुग्मित (3d में 5, 4s में 1)।
\textbf{17.} \ce{Cu} $[\ce{Ar}]\,3\text{d}^{10} 4\text{s}^1$, \ce{Cu+}
$[\ce{Ar}]\,3\text{d}^{10}$, \ce{Cu^{2+}} $[\ce{Ar}]\,3\text{d}^9$।
% ledger: gs:Cu, gs:Cu+, gs:Cu2+
\textbf{18.} \ce{Cu^{2+}}, एक \omterm{def:b1:quantum-numbers:unpaired}{अयुग्मित इलेक्ट्रॉन} के साथ; \ce{Cu+} में कोई नहीं।
\textbf{19.} \ce{Sc^{3+}} $[\ce{Ar}]$: 0; \ce{Ti^{3+}} $[\ce{Ar}]\,3\text{d}^1$:
1; \ce{Mn^{2+}} $[\ce{Ar}]\,3\text{d}^5$: 5। सबसे अधिक \ce{Mn^{2+}} में हैं।
\textbf{20.} हर कक्षक में तीन इलेक्ट्रॉन; हर \omterm{def:b1:quantum-numbers:orbital}{उपकोश} में $3(2l + 1)$;
हर कोश में $3n^2$।
\textbf{21.} 1s: 3; 2s: 3; 2p: 9; 3s: 3; 3p: 9।
\textbf{22.} पहला आवर्त 1s भरने पर बंद होता है: $Z = 3$।
\textbf{23.} $Z = 4$, उत्कृष्ट क्रोड के बाद 2s में एक इलेक्ट्रॉन।
\textbf{24.} 2s और 2p: $3 + 9 = 12$ तत्व।
\textbf{25.} $3 + 12 = 15$: तीन-चक्रण संसार की दूसरी उत्कृष्ट गैस का
परमाणु क्रमांक \textbf{$Z = 15$} है।
\end{solution}
